% Section 4: Formal Definitions
% Target length: 450-700 words (~0.75-1 page). Keep tight -- this is a definitions
% section, not prose-heavy.
% Status: Definitions and Proposition 1 are drafted at working-notes level below;
%         needs LaTeX theorem-environment polish and a proof or proof sketch.

\section{Formal Definitions}
\label{sec:formal-definitions}

\subsection{Learning Graph}

\textbf{Definition 1 (Learning Graph).}
A \emph{Learning Graph} is a monopartite directed graph $G = (C, D)$ where $C$
is a set of concept nodes and $D \subseteq C \times C$ is a set of dependency
edges. An edge $(a, b) \in D$ means concept $a$ depends on (requires) concept
$b$. $G$ is constrained to be acyclic: a circular dependency would make the
prerequisite relation unsatisfiable. $G$ is therefore a DAG.

% TODO: cite this repo's own learning-graph-schema.json / documentation as the
% concrete implementation of this definition (nodes have id/label/group/shape;
% edges have from/to). Confirm edge direction convention matches (from=dependent,
% to=prerequisite) -- verified empirically on the Algebra I graph in Section 6.

\textbf{Definition 2 (in-degree).}
$\mathrm{indeg}(b) = |\{a : (a,b) \in D\}|$, the number of concepts that
\emph{directly} depend on $b$.

\subsection{Concept Impact Score (CIS)}

% TODO: motivate in 1-2 sentences why in-degree is insufficient before the
% definition (transitive/indirect dependents undercounted -- see Algebra I
% example data in Section 6, Table 1).

\textbf{Definition 3 (Concept Impact Score).}
\[
\mathrm{CIS}(b) \;=\; \frac{1-d}{|C|} \;+\; d \sum_{a\,:\,(a,b)\in D} \frac{\mathrm{CIS}(a)}{\mathrm{outdeg}(a)}
\]
This is standard PageRank applied directly to $G$. Note that the dependency
edge convention ($a \to b$ means ``$a$ depends on $b$'') already matches the
citation convention (citing paper $\to$ cited paper), so unlike some
applications no graph transpose is required.

\textbf{Proposition 1 (Exact single-pass computation on a DAG).}
Because a Learning Graph is acyclic, CIS admits an exact closed-form solution
computable in a single pass, without iteration and without a damping factor:
process concepts in topological order (each concept only after every concept
that depends on it has already been scored), and
\[
\mathrm{CIS}(b) = 1 + \sum_{a\,:\,(a,b)\in D} \mathrm{CIS}(a)
\]
falls out directly, in $O(|C|+|D|)$ time.

% TODO: write proof sketch. Outline:
%  1. G is a DAG => has >=1 valid topological order (standard result).
%  2. Process nodes in that order; by definition every predecessor a with
%     (a,b) in D is processed before b, so CIS(a) is already finalized.
%  3. No back-edges exist (acyclicity), so no node is revisited => the single
%     pass yields the unique fixed point that damped iterative PageRank would
%     otherwise converge to in the limit.
%  4. Contrast with web-graph PageRank, where damping (d < 1) exists SPECIFICALLY
%     to guarantee convergence in the presence of cycles and dangling nodes --
%     conditions that cannot occur in a Learning Graph by Definition 1.
%  Relate to TotalRank \cite{boldi2005} in Related Work, but note this is a
%  domain-driven simplification, not a general-graph technique.

\textbf{Design decision: global normalization (resolved).}
Normalization scope was decided as \textbf{global}: a concept's elaboration
tier is determined relative to the \emph{entire} book's concept set $C$, not
just the concepts co-occurring in its chapter. Rationale: ``impact on the
entire textbook'' is a book-level claim by definition
(\S\ref{sec:problem-statement}), and local normalization produced degenerate
results in practice -- see \S\ref{sec:system}, where a chapter-local
normalization forced 16 of 17 Chapter~1 concepts into the lowest tier purely
because that chapter happens to contain the single globally-highest-CIS
concept in the book.

\textbf{Definition 4 (Elaboration Score).}
The first global-normalization attempt used \emph{population percentile
rank} -- a concept's tier was set by what fraction of all $|C|$ concepts it
outranked. This failed a second validation check (Chapter 12; see
\S\ref{sec:system}): because CIS has a hard floor of 1 (every concept
trivially satisfies $\mathrm{CIS}(b) \geq 1$ from its own base case) and a
large share of concepts in a real learning graph are true leaves (103 of 200
in the Algebra~I graph, 51.5\%), the population 35th-percentile cut point
sits \emph{at} the CIS floor -- making the lowest tier structurally
unreachable regardless of how unimportant a concept actually is. We instead
define the \emph{Elaboration Score} as a normalized log-scaled value, not a
population rank:
\[
E(c) \;=\; \frac{\log(\mathrm{CIS}(c) + 1)}{\log(\mathrm{CIS}_{\max} + 1)} \in [0, 1]
\]
where $\mathrm{CIS}_{\max}$ is the maximum CIS over the whole book. Tied
concepts at the CIS floor now separate cleanly from true hubs ($E=0$ is the
floor, $E=1$ is the book's single most-depended-upon concept), rather than
being absorbed into the same percentile band as everything else near the
median. This is also consistent with the word-budget weighting formula
(\S\ref{sec:problem-statement}), which already used $\log(\mathrm{CIS}+1)$ as
its weight -- tiering and word-budget weighting now share one normalized
scale instead of two different ones (population rank for tiers, log-value for
words).

\textbf{Tier assignment (revised):} Tier A if $E(c) \geq 0.5$; Tier B if
$0.2 \leq E(c) < 0.5$; Tier C if $E(c) < 0.2$. Validated on two chapters with
opposite expected profiles (\S\ref{sec:system}, Tables~1-2): Chapter~1
(foundations) yields 8 A / 6 B / 3 C -- still foundation-skewed but no longer
degenerate; Chapter~12 (a late, specialized chapter of mostly terminal
concepts) yields 0 A / 2 B / 8 C -- the mirror-image profile a specialized
chapter should have. This cross-chapter contrast is the validation check the
first (percentile-rank) scheme failed.

% TODO: remaining open design choice -- undiluted sum-of-descendants CIS
% (Proposition 1 form, used above) vs. normalized/diluted PageRank-style CIS
% (sums to 1 across C). Does not affect the global-vs-local decision, and does
% not affect E(c)'s qualitative behavior (log of a rescaled quantity is still
% monotonic), but changes the exact E(c) values -- confirm this doesn't shift
% any concept across a tier boundary before finalizing.
%
% TODO: the 0.5/0.2 cut points for E(c) are themselves a fresh, only lightly
% validated choice (two chapters). Treat as a starting point, not a settled
% constant -- revisit once more chapters are run through this table
% (\S\ref{sec:evaluation}).

% TODO: add Figure 1 -- small toy graph (6-8 nodes) illustrating a case where
% in-degree ranking and CIS ranking disagree, annotated with both scores.
% Real candidate data already computed: data/algebra1-concept-impact.csv
% (Constant: indeg=1 rank #53, CIS=238 rank #4; Coefficient: indeg=2 rank #35,
% CIS=239 rank #3).

\textbf{Methodological note worth keeping in the paper.} The percentile-rank
failure above is a useful cautionary example in its own right: a recursive
importance measure with a hard floor and a heavy-tailed distribution (as CIS
is on a real learning graph -- most concepts are leaves) should not be
discretized by population rank without checking for exactly this
floor-saturation failure mode. This is a generalizable point beyond this
paper's specific metric and may be worth a short remark in
\S\ref{sec:discussion}.
